% Ad Familiares 13.32 --- headnote
% ad-familiares-13-32, 46 BC, Rome

\textit{Cicero to Manius Acilius Glabrio, proconsul of
Sicily, written from Rome in 46 BC (the manuscript
dateline: \textit{Scr. Romae, ut videtur a. 708 (46)}).
This is one of a cluster of recommendation letters
(Fam.\ 13.30--39) addressed to Acilius on behalf of
clients with affairs in his province. The beneficiaries
here are the brothers Marcus and Gaius Clodius
Archagathus and Philo of Halaesa, a prosperous and
distinguished Sicilian town of which their family was
evidently a leading house; Cicero's tie to them is one
of long standing, formed of hospitality, mutual
obligation, and personal regard.}

\textit{The interest of the piece is the open
acknowledgement of the rhetorical embarrassment of the
genre. Cicero, sending Acilius a sheaf of these notes,
is aware that by commending so many people in the
strongest terms he risks flattening the currency of
recommendation altogether --- and he says so, in a brief
and slightly rueful aside, before pressing on with the
request anyway. The candour is itself a Ciceronian
gesture: the writer who knows that recommendations are
counters in a patronage economy, and trusts his
correspondent to know it too, can afford to admit the
fact without ceasing to ask.}
